%%=============================================================================
%% CAPITULO 1 -- INTRODUCAO
%% Versao dummy: resumo do capitulo 01-introducao da dissertacao real.
%%=============================================================================

\section{Introdução}

O tema desta pesquisa é o uso de uma arquitetura de Inteligência Artificial
(IA) Generativa com recuperação de conhecimento (\textit{Retrieval-Augmented
Generation}, RAG) e cálculo determinístico para democratizar a assessoria em
renda fixa em bancos públicos. O número de investidores pessoa física cresceu,
mas a orientação financeira personalizada e tecnicamente precisa continua
acessível sobretudo a clientes de alta renda. Os robôs-investidores
tradicionais limitam-se a questionários estáticos de \textit{suitability}, e os
modelos de linguagem de larga escala (LLMs) usados sozinhos estão sujeitos a
alucinações matemáticas, risco inaceitável em ambiente regulado. Os bancos
públicos, federais e regionais como o Banco de Brasília (BRB), têm mandato
social e ampla base de correntistas de perfil popular, mas não conseguem
ampliar o quadro de especialistas na mesma proporção do crescimento da base de
clientes.

Diante desse quadro, a pergunta de pesquisa é a seguinte: uma arquitetura de
IA Generativa baseada em agente único que, por meio de \textit{function
calling}, separa o cálculo quantitativo em Python da geração de linguagem por
LLM e recupera normativos vigentes de uma base de conhecimento regulatória
(RAG) é capaz de oferecer orientação em investimentos de renda fixa para
clientes pessoa física de bancos públicos com precisão numérica verificável e
conformidade à Resolução CVM nº~30/2021, com adequação técnica comprovada por
comparação automatizada contra padrões de referência verificáveis? O objetivo
geral é desenvolver e validar essa arquitetura, com foco em acessibilidade,
precisão numérica verificável e conformidade regulatória, contribuindo para a
literatura de IA aplicada à Administração Pública e para a democratização do
acesso à assessoria financeira de qualidade.
